\section{Introduction}
\label{sec:introduction}

As the threat posed by quantum computers becomes increasingly realistic, the migration from conventional public-key cryptography to post-quantum cryptography (PQC) is rapidly progressing. For this transition to be deployed in real-world systems, not only security but also implementation efficiency is required.

\textbf{AIMer} is a digital signature scheme selected as a final algorithm in the Korean Post-Quantum Cryptography (KpqC) competition. Many PQC signature schemes rely on mathematical structures such as lattices or multivariate quadratic equations. In contrast, AIMer derives its security solely from a symmetric-key one-way function and an MPC-in-the-Head (MPCitH) zero-knowledge proof~\cite{kim2023aimer}. The primary challenge of such a symmetric-primitive-based design is its computational cost. However, the computation is inherently parallel. In particular, the signer repeatedly performs \emph{the same binary-field arithmetic} for a large number of virtual MPC parties. The entire process is tied together by \emph{a large number of Keccak permutations}. Both components are ideal targets for wide vector instructions (SIMD). This observation motivates the work presented in this paper.

Despite this opportunity, the AIMer implementation released in January 2026 provides only a portable C reference implementation~\cite{choaimer}. KpqClean, which collects implementations of KpqC candidate algorithms, includes AVX2-based optimizations. However, it does not yet exploit AVX-512, a 512-bit SIMD instruction set available on modern processors. This instruction set includes VPCLMULQDQ for binary polynomial multiplication and VPTERNLOGQ for three-input bitwise logic. In this work, we integrate both the reference and AVX2 implementations into the liboqs (Open Quantum Safe) framework. Liboqs is an open-source library that provides a common interface for a wide range of post-quantum cryptographic algorithms. We use these implementations as baselines and extend them with AVX-512 optimizations.

To address the issues outlined above, this paper makes the following contributions:

\begin{enumerate}[label=\arabic*., leftmargin=*]
\item We present the first comprehensive AVX-512 implementation covering all six AIMer variants. The implementation includes a VPCLMULQDQ-based four-party field layer, a VPTERNLOGQ-based linear layer, and an AVX-512VL-based Keccak implementation. Compared with the AVX2 implementation, it achieves a speedup of $1.60$--$1.84\times$.

\item We analyze the additional speedup obtained when moving from AVX2
to AVX-512. The AVX2 baseline already vectorizes binary-field arithmetic.
In the 128-bit and 192-bit variants, Keccak accounts for approximately
86\% of this additional improvement. These ratios depend on the selected baseline.

\item We analyze the variation in AVX-512 efficiency across AIMer variants from the perspectives of field width and the number of MPC parties. Elements in $\gftwo{128}$ and $\gftwo{256}$ can be packed
efficiently into 512-bit registers, resulting in higher gains. In contrast, $\gftwo{192}$ incurs losses due to zero padding. Furthermore, the small variants benefit more from four-party packing because the packing operation is repeated more frequently as the number of parties increases.
\end{enumerate}

The remainder of this paper is organized as follows. Section~\ref{sec:bg} reviews the background of the AIMer signature scheme, the AIM2 one-way function, and the MPCitH proof system. Section~\ref{relatedwork} surveys prior work on MPCitH signatures and SIMD-based cryptographic optimization. Section~\ref{sec:impl} describes our AVX-512 implementation, including VPCLMULQDQ-based four-party field arithmetic, a VPTERNLOGQ-based linear layer, and AVX-512VL-accelerated Keccak. Section~\ref{sec:eval} quantitatively evaluates performance relative to AVX2 and analyzes the observed improvement by decomposing it into field-arithmetic and Keccak contributions. Finally, Section~\ref{sec:conc} concludes the paper and discusses limitations and future research directions.
